\section{Track progress across short file operations}
A successful open returns a descriptor that belongs to that open operation. Save it;
do not assume it is always three. For a read request, distinguish a positive count,
zero at end of file, and a negative error. A positive result can be smaller than the
request. Advance the destination pointer and accumulated count by the result, then
continue. Treating the first short read as EOF can lose the rest of a valid image.

When the input arena is full, one more read into a separate one-byte private probe can
distinguish exact-fit EOF from an overlarge file. Do not put that probe beyond the
caller's arena. Check an error from the probe separately from EOF. Every successful
open needs a close on success, malformed input, capacity failure and I/O error paths.

Writing uses the same progress invariant. Keep a current source pointer and a remaining
byte count. After a positive result, advance the pointer by that count and subtract it
from the remainder. Retry the remainder, not the original request. A negative result
is an error. A zero result for a nonempty request cannot make progress and must also
terminate with an error rather than loop forever.

\section{Account for the current syscall's read footprint}
The current MARS write service reads one byte after the requested interval before passing
the buffer to its file environment. This is an existing educational ABI behavior, verified
by the core boundary test. It also occurs when the file environment later accepts a
short prefix. The file count does not include that extra byte.

Provide a readable guard byte after the complete serialized file buffer. It is storage,
not an additional PPM sample and not a terminator written to the output file. On the
last request, the guard supplies the service's extra read. On earlier short requests,
the next serialized byte lies within the still-readable buffer. A trace should account
for this source read while continuing to forbid writes outside the actual file interval.

The grading fixtures now distinguish declared input and output paths. A write can be
shortened or failed at a chosen byte count, and the resulting bytes can be inspected
without opening a host file. The file operations themselves remain observable, including
the supplied descriptor, requested count, returned count and close. An image that merely
looks correct cannot establish these resource and error-path properties.
